\chapter{Experimental Design}\label{ch:experimental-design}

This chapter fixes the configuration under which the method of
Chapter~\ref{ch:methodology} was instantiated: the corpus and its
split, the rosters, the conditions, the pre-registrations with what
each of them fixed in advance, and the statistical protocol.
Taxonomy-tuning mechanics and the construction ablation matrix are in
Appendix~\ref{app:tuning-protocol}; rubric-generation mechanics are in
Appendix~\ref{app:judge-generation}.

\section{Research Questions, Conditions, and Metrics}\label{sec:exp-rqmap}

Table~\ref{tab:rq-map} fixes which condition answers which research
question, with which metric, and which link of
Section~\ref{sec:research-questions} it lands on.

\begin{table}[h]
    \centering
    \caption{Research question to condition, metric, and link.}
    \label{tab:rq-map}
    \begin{tabular}{p{0.9cm}p{4.4cm}p{3.0cm}p{4.4cm}}
        \toprule
        RQ & Question & Conditions & Primary metric \\
        \midrule
        RQ1 & Do answer-key-blind, domain-conditioned pairwise
        judgments recover per-domain rankings, and what decides the
        verdicts?
        & C1 (main), with C3 as the rubric control
        & Per-verdict agreement with the MMLU-Pro key; Spearman $\rho$
        against ground-truth accuracy, reported under both tie
        conventions \\
        \addlinespace
        RQ2 & Does adapted, cell-scoped grouping enable better judging
        than a generic partition-free rubric on the same questions?
        & C1 vs.\ C5, paired on an identical question set
        & Spearman $\rho$ per domain (partition claim); per-verdict key
        agreement (judge-side check) \\
        \bottomrule
    \end{tabular}
\end{table}

The scope is a \textbf{selected-domain study}: the arena evaluates one
or two ex-ante-selected, high-support domains at high power rather than
all 14 at matched power, and the selection rule is a compute-budget
trade-off stated in Section~\ref{sec:exp-datasets}.

\section{Implementation Environment}\label{sec:exp-env}

The pipeline is implemented in Kotlin on Spring Boot, backed by SQLite
in WAL mode (Section~\ref{sec:arch-infra}). Embeddings are produced by
Qwen3-Embedding served locally; the judge is a single large
instruction-tuned model accessed through an Azure OpenAI-compatible
client, and the same judge model is used for rubric induction and for
every arena condition.

\textbf{Every construction and arena result in this thesis is at seed
42.} No result is reported as a mean over multiple seeds, and no
three-seed convention is in force. Where seed dispersion was measured
--- in the acceptance-gate analysis of
Section~\ref{sec:res-construction} --- it is reported as a separate,
dated ablation over five seeds and not as an error bar on a headline
number.
The canonical configuration values are tabulated with their provenance
in Appendix~\ref{app:tuning-protocol}
(Table~\ref{tab:hyperparams}). Per-run identifiers, seeds, and realised
call counts are recorded in the manifest emitted with each run.

\todo[inline]{Confirm the judge model identity against the run manifest
for the reported runs, and make the roster table state it. Two sources
in this project disagree, and Section~\ref{sec:exp-rosters}'s
``the judge appears in no roster, so there is no self-preference
confound'' depends on resolving it.}

\section{Corpus, Split, and Information Separation}\label{sec:exp-datasets}

The corpus is MMLU-Pro~\parencite{wang2024mmlupro}: 12{,}000 questions
across 14 domains with human-verified reference answers, used in two
roles. It supplies the ground-truth domain labels that seed the 14
anchor nodes, and it supplies the per-domain accuracy scores used as
the post-hoc validation oracle (Section~\ref{sec:meth-validation}).
All 14 categories seed anchors, including the semantically
heterogeneous \textit{Other} category, which is retained as a geometric
anchor --- its queries remain routable mass and may migrate toward
coherent anchors --- but is not an arena selection candidate.
AG News~\parencite{zhang2015agnews} serves only as a pipeline sanity
check on an independently labelled corpus; it seeds no anchors and
enters no arena.

Taxonomy construction runs on all 14 domains simultaneously, which is
what allows queries to migrate between anchors and so makes the adapted
partition differ from the canonical one at all.
The arena evaluates a roster on one or two ex-ante-selected domains,
chosen before observing any arena result under a fixed rule: high
held-out query support; high geometric coherence relative to siblings;
non-trivial adapted-versus-canonical migration volume; and
non-degenerate model performance variance under existing MMLU-Pro
accuracy. Math was the primary candidate under that rule. The offline
domain screen of Section~\ref{sec:res-law} was added later and is a
stronger instrument than criterion~(iv), because it asks directly
whether a domain's models reorder relative to their global ordering;
it selects law, and it demotes Math to a null arm.

\subsection{The Split}

The corpus is partitioned into three disjoint, stratified subsets:
\textbf{construction} (taxonomy induction and rubric generation),
\textbf{tuning-validation} (tuning gates only), and
\textbf{arena-test} (held-out routing and judging only), stratified by
MMLU-Pro sub-topic tag and fixed across seeds.

The intended proportions are 60/15/25, and the implementation does not
match them. The runner's default is a $70/30$ split, and the
tuning-validation partition is enforced by scoring configurations on
construction-internal metrics rather than as an explicitly materialised
subsample. No arena-test query is seen by any tuning gate under either
realisation, so the no-leakage property on which the answer-key-blind
claim rests holds in both; but the split as run is $70/30$ and is
described that way wherever a query count is quoted.

One further detail belongs here because it affects what a
``construction-set'' claim means: the rubric-induction pool is
additionally thinned by roughly a fifth by a hash filter, so a leaf's
rubric is induced from a sample of its construction queries rather than
all of them.\footnote{And about $4.2\%$ of MMLU-Pro rows --- some 500
of 12{,}000 --- carry no link into the evaluation tables; they are
believed to be sentinel identifiers, but that has not been confirmed,
and they are recorded here as unattributed rather than explained.}

The arena-test fraction targets approximately 10 held-out queries at
the smallest leaf under a selected domain. The binding constraint is
the number of model pairs per leaf, $\binom{11}{2} = 55$ at the
corrected roster.

\subsection{Information Separation}\label{sec:exp-info-separation}

Table~\ref{tab:info-separation} makes the answer-key-blind design
concrete: at every pipeline phase, only the data marked \checkmark{} is
accessible to the corresponding component.

\begin{table}[h]
    \centering
    \caption{Data access by pipeline phase.}
    \label{tab:info-separation}
    \begin{tabular}{p{3.6cm}ccccc}
        \toprule
        Phase &
        $\mathcal{Q}_{\text{con}}$ text &
        $\mathcal{Q}_{\text{con}}$ labels/answers &
        $\mathcal{Q}_{\text{eval}}$ text &
        $\mathcal{Q}_{\text{eval}}$ labels/answers \\
        \midrule
        Taxonomy construction & \checkmark & \checkmark & --- & --- \\
        Rubric generation & \checkmark & \checkmark & --- & --- \\
        Held-out routing & --- & --- & \checkmark & --- \\
        Arena evaluation (judging) & --- & --- & \checkmark & --- \\
        Post-hoc validation & --- & --- & \checkmark & \checkmark \\
        \bottomrule
    \end{tabular}
\end{table}

The evaluation set's labels and reference answers are accessed only in
the final, post-hoc validation phase, after every verdict has been
recorded. The rubric-induction filter that enforces the second row is
audited by a build-failing assertion
(Section~\ref{sec:arch-judge}), which is the strongest guarantee in the
system.

It guarantees exactly what it says and no more. The judge does not
receive the answer key; it does receive the multiple-choice options and
responses that usually state their own answer. ``Answer-key-blind'' is
not ``answer-blind'', and Chapter~\ref{ch:results} turns on the
difference.

\section{Rosters and the Exclusion Policy}\label{sec:exp-arena}\label{sec:exp-rosters}

Three rosters appear in this thesis. They are not the same roster, and
saying so is more useful than smoothing them into one.

\textbf{The pilot roster (8 models)} is the one Run~A used. It is
reported with its defect rather than repaired retroactively: four of
its eight models had their response text stored as the empty string
across all rows, and the resulting format artifact decided half of that
run's comparisons (Section~\ref{sec:res-capture-gap}).

\begin{table}[h]
\centering
\caption{The pilot roster (Run~A), with stored response format and
ground-truth accuracy on the 241 judged Math questions. The two format
groups do not overlap in accuracy, so trace presence and capability are
collinear in this roster.}
\label{tab:model-roster-pilot}
\begin{tabular}{llr}
\toprule
Model & Stored response & Math accuracy \\
\midrule
gpt-4o-2024-08-06 & reasoning, median 1211 ch & 78.4\% \\
claude-3-5-sonnet-20241022 & reasoning, median 630 ch & 75.9\% \\
deepseek-chat-v2\_5 & reasoning, median 1038 ch & 71.0\% \\
claude-3-5-haiku-20241022 & reasoning, median 664 ch & 57.7\% \\
\addlinespace
Meta-Llama-3\_1-70B-Instruct & \textbf{empty; stub synthesised} & 53.1\% \\
Qwen1.5-72B-Chat & \textbf{empty; stub synthesised} & 36.5\% \\
Llama-2-70b-hf & \textbf{empty; stub synthesised} & 13.3\% \\
Llama-2-13b-hf & \textbf{empty; stub synthesised} & 7.5\% \\
\bottomrule
\end{tabular}
\end{table}

\textbf{The length-matched band (12 models)} is the roster of Run~B, on
the repaired corpus. It is selected so that response length and
accuracy are close to uncorrelated within the roster
($r = +0.286$, against $+0.597$ across all 36 valid models), which
controls the format artifact by construction rather than qualifying it
afterwards. All its members produce substantive responses; answer-only
systems are out of scope for this evaluation design, because a rubric
has nothing to grade in a response with no content.

\textbf{The corrected band (11 models)} is what the exclusion policy
yields at full law coverage: the length-matched band less the one
member that emits a bare answer on $35.4\%$ of items. It is the roster
of the offline domain screen, of the reliability-constant fit, and of
Run~C. Post-repair, its members are at least $99.8\%$ substantive.

\paragraph{Exclusion is enforced, and it is a measurement rather than a
housekeeping decision.}
Ten models are excluded at roster load by a hard-coded list that
raises, each entry carrying its reason: one identifier-space mismatch,
six archives with no stored reasoning text, and three mixed-format
models. The mixed-format threshold is format \emph{consistency}, not
text presence: a model that emits a bare answer on a third of items
reintroduces the format preference inside its own comparisons. The
three excluded on that ground emit bare answers on $68.4\%$, $35.4\%$
and $27.6\%$ of items respectively.

The honest statement of the policy is therefore:
\emph{response text was not persisted by the generation pipeline for
several models, and three further models emit a bare answer on more
than a quarter of items; all are excluded, and the format artifact is
reported as the measurement that justifies the exclusion.}

One note belongs with it: six of the ten exclusions are for exactly
the condition the corpus validator's trace-presence check measures and
declines to act on, because that check has no failure branch. The
policy is a ratchet on known failures, not a threshold, and it does
not generalise to any new model.\footnote{And one disambiguation: the
roster member \texttt{Meta-Llama-3\_1-70B-Instruct} is not the
excluded \texttt{Meta-Llama-3-70B-Instruct}. They are different
models, the roster's is clean, and the excluded one sits in a
different question-identifier space and is what collapsed law to 20
questions in Section~\ref{sec:res-law}.}
The full roster tables are in Appendix~\ref{app:models}.

\section{Conditions}\label{sec:exp-arena-conditions}

\begin{table}[h]
    \centering
    \caption{Arena conditions. All conditions share the frozen
    snapshot and the model roster of their run; only the listed factor
    varies.}
    \label{tab:arena-conditions}
    \begin{tabular}{p{0.9cm}p{2.7cm}p{4.4cm}p{4.0cm}}
        \toprule
        ID & Condition & Delta from C1 & What it controls \\
        \midrule
        C1 & Main
        & --- (adapted partition, cell-scoped reference-informed rubric)
        & reference \\
        \addlinespace
        C3 & Generic judge
        & Same cells, same triples; the judge prompt is a shared
        structural template with no cell-specific rubric
        & rubric specificity, partition held fixed \\
        \addlinespace
        C5 & No partition
        & Same questions; judged under a generic MT-Bench-style
        pairwise prompt, with Bradley--Terry fitted directly at the
        domain
        & \textbf{the partition itself} \\
        \bottomrule
    \end{tabular}
\end{table}

C3 and C5 control different things, and the difference is why both
exist. C3 replays C1's leaf-scoped triples, so it \emph{shares the
partition it would have to be a control for}; it can isolate the rubric
and nothing else. C5 removes the partition. It is the no-taxonomy
baseline.\footnote{A fourth condition, re-aggregating C1's verdict set
by editorial label, was specified in an earlier design document and is
cut: it shares its entire verdict set with C1 and therefore cannot
isolate anything, and the question it was built to answer --- the
per-cell adapted-versus-canonical rank difference --- was refuted at
its premise under pre-registration before the arena ran
(Section~\ref{sec:res-granularity}).}

\paragraph{Pairing, and what may be compared across formats.}
Both arms of the C1-versus-C5 comparison judge an identical question
set. Unlike the original MT-Bench protocol, which samples
independently because it has no partition to compare against, here the
partition is the treatment, so pairing is what isolates it; independent
sampling would add variance orthogonal to the thing under test.
The two arms parse different verdict formats --- a structured schema
against \texttt{[[A]]}/\texttt{[[B]]}/\texttt{[[C]]} markers --- with
three consequences. Winners, the Bradley--Terry ordering derived from
them, and tie counts \emph{are} comparable, provided ties are defined
identically on both sides. Judge confidence is \emph{not}: the generic
arm emits a flat constant. Parse-failure rates are not comparable
either, and the failure \emph{modes} differ rather than merely the
rates, because a malformed structured response fails loudly while a
generic response with no marker returns an invalid verdict silently.
Operationally this forces one thing that is easy to miss: the
confidence gate must be disabled in both arms, since at a flat
confidence it would admit every generic verdict while filtering the
main arm's, and the arms would no longer be paired where it counts.

\paragraph{And a precondition that came out of C3.}
Before spending a judging budget on any A/B arm, \textbf{measure
verdict agreement first}. It is one query against the verdict store. C1
and C3 turned out to agree on $99.4\%$ of winners
(Section~\ref{sec:res-generic-judge}), which means no ranking
difference between them was detectable at any sample size --- a fact
obtainable before the second arm was run rather than after.

\section{Pre-Registration}\label{sec:exp-prereg}\label{sec:exp-baselines}

Four decision rules were fixed in dated documents before the data they
govern existed. They are reproduced in full in
Appendix~\ref{app:tuning-protocol}; what each fixed is summarised here,
because Chapter~\ref{ch:results} calls results ``pre-registered'' and a
reader is entitled to check the claim.

\textbf{Granularity and discriminative power.} Three outcomes were
named in advance, and the primary statistic was fixed as the
\emph{observed} correlation rather than the disattenuated one. That
rule is the only reason the conclusion is readable: the disattenuated
column moves in the opposite direction and clips past $1.0$, so leading
with it would have produced a conclusion from the arithmetic of the
correction rather than from the data.

\textbf{Rubric specificity.} Four directional predictions, with the
confound direction named \emph{per measure} --- including one measure
whose vocabulary-size confound could have driven the null arm
\emph{higher} --- plus within-arm decision thresholds.

\textbf{Arena launch cost.} Cost bands with the instruction that no
band is to be renegotiated after the number is seen, and a void
condition: if the measured per-cell identification cost came back
constant across cells, the scheduler would still be capping and the
measurement would report nothing.

\textbf{The no-partition baseline.} The paired design; the primary and
secondary outcomes; the decisive difference of two Spearman grid steps
on Math; four void conditions (identical question sets, confidence gate
off in both arms, identical samplers, identically defined ties); and
--- recorded before the run --- the honest prior that with options
present the comparison would likely come back null, together with the
reason.

The baseline conditions for taxonomy quality specified in an earlier
design --- flat $k$-means, agglomerative clustering, and a random
taxonomy null, scored by a family of reference-free hierarchy metrics
--- are not reported. The supervised hierarchical-classification
baselines were removed as unimplemented; no gold taxonomy exists for
MMLU-Pro, so reference-dependent metrics were dropped rather than
computed against an ad-hoc structure; and the remaining runs were never
executed. None of those metrics bears on any of the three links of
Section~\ref{sec:research-questions}, and the construction evidence
reported instead is the fixed-point certificate, the leaf-size
distribution, and the separation nulls of
Section~\ref{sec:arch-requirements}.

\todo[inline]{The rubric-specificity pre-registration is stale as
written: it still records ``treatment arm only, no null, this is not
yet evidence''. The null arm has since been run and its results were
never written back. Amend the document, with the amendment dated, before
reproducing it in the appendix.}

\section{Metrics and Statistical Protocol}\label{sec:exp-metrics}

\textbf{The primary judge-side metric is per-verdict agreement with the
MMLU-Pro answer key}, on the subset where exactly one response is
correct and the judge did not tie. It is a proportion over roughly a
thousand verdicts per condition, with a standard error near $1.5\%$,
and it resolves differences of about four points.

\textbf{Spearman $\rho$ against ground-truth accuracy is the partition
metric, and it is heavily quantised.} At $M$ models the grid spacing is
$6/(M(M^2-1))$: $0.0119$ at $M = 8$, $0.0060$ at $M = 10$, $0.0035$ at
$M = 12$. A single Spearman value therefore consumes a full arena run to
produce one number resolvable only to an adjacent transposition, which
is why it is the secondary statistic everywhere except the partition
claim, where it is the thing being claimed about.

\textbf{Bradley--Terry $\theta$, never win-rate.} The scheduler
over-samples uncertain pairs, so mid-ranked models face harder average
opposition; win-rate ignores opponent strength and is biased by this,
and $\theta$ is not. The measured cost of the difference --- the
$5/0$-versus-$5/2/4$ sign-test contrast --- is shown in
Section~\ref{sec:meth-scheduling}.

\textbf{Every $\rho$ is reported under both tie conventions} ---
ties half-weighted and tied comparisons dropped, in that order --- per
rule~D5 of Section~\ref{sec:meth-discipline}; the measured instability
that forces the rule is Section~\ref{sec:res-tie-policy}.

\textbf{Disattenuation is two-sided.} Both sides of
$\rho(\text{arena ranking}, \text{ground-truth ranking})$ are
finite-sample per-cell rankings, so the correction is $\rho_{\text{obs}}
/ r$ with $r = n/(n+c)$, not $\rho_{\text{obs}} / \sqrt{r}$; the
one-sided form applies only when one measure is perfectly reliable. The
corrected constant is $c = 2.52$ at the 11-model band, and three
caveats travel with it wherever it is used: it is a property of the
roster and must be re-fitted if the roster changes; the interval
originally quoted for it resamples three fitted points and is not a
confidence interval; and it was fitted on ground-truth rankings, which
are noise-free per question, so it is a \emph{ceiling} on arena
reliability rather than a prediction of it
(Section~\ref{sec:disc-debt}).

\textbf{Bootstrap resampling counts differ by path and are stated where
used.} The validation path uses 2{,}000 resamples; the ranking path
uses 50; the in-run trajectory uses 1. A blanket claim of 2{,}000
resamples everywhere would be wrong, and figures derived from the
ranking path are quoted with their own count.

\textbf{No pass/fail threshold on $\rho$.} An earlier design
pre-specified $\rho > 0.7$ as a criterion. It is retained only as an
interpretive marker of strong agreement; results below it are reported
in full and scoped to the domains they were measured on.
